% Generated from OC_CORE_1_3_SCIENCE_SPOT_latest.json
\subsection{Компетенция закона ревизий}
\begin{center}
\begin{tabular}{@{}L{0.08\textwidth}L{0.17\textwidth}L{0.23\textwidth}L{0.24\textwidth}L{0.22\textwidth}@{}}
\toprule
Приоритет & Слой & Правило ревизии & Триггер конфликта & Текущий статус \\
\midrule
1 & \occode{EMPIRICAL_BRIDGES_AND_REPLAY_BINDINGS} & \occode{REVISE_EMPIRICAL_BRIDGES_FIRST} & \occode{OFFICIAL_OR_INSTITUTE_RUN_REPLAY_BREACH} & \occode{STABLE} \\
2 & \occode{PARAMETERIZATION_AND_NUMERICAL_PACKETS} & \occode{REVISE_PARAMETERIZATIONS_BEFORE_THEOREM_SCOPE} & \occode{RESIDUAL_OR_CALIBRATION_BREACH_PERSISTS_AFTER_BRIDGE_REPAIR} & \occode{STABLE} \\
3 & \occode{STANDALONE_DOMAIN_PACKETS_AND_DERIVED_THEOREMS} & \occode{REVISE_STANDALONE_DOMAIN_PACKETS_BEFORE_ROOT_KERNEL} & \occode{MULTI_BENCHMARK_FAILURE_PERSISTS_AFTER_PARAMETERIZATION_REPAIR} & \occode{REAUDIT_REQUIRED_FOR_BRIDGE_ONLY_DOMAINS} \\
4 & \occode{ROOT_KERNEL} & \occode{ROOT_KERNEL_REVISION_ALLOWED_ONLY_AFTER_ACCUMULATED_CROSS_DOMAIN_FAILURE} & \occode{ACCUMULATED_CROSS_DOMAIN_FAILURE_EXCEEDS_DECLARED_THRESHOLD} & \occode{LOCKED_PENDING_ACCUMULATED_CROSS_DOMAIN_FAILURE} \\
\bottomrule
\end{tabular}
\end{center}

\subsection{Доказательные обязательства, несущие нагрузку}
\begin{center}
\begin{tabular}{@{}L{0.18\textwidth}L{0.28\textwidth}L{0.22\textwidth}L{0.26\textwidth}@{}}
\toprule
ID обязанности & Формулировка & Исключённые зависимости & Условие фальсификатора \\
\midrule
\occode{OBLIGATION::SOURCE_THEOREM_3} & Показать, что Axiom 3.2 вместе с Definitions 12.1, 12.3 и 12.4 строго влечёт Source Theorem 3. & Исключённое семейство lift-теорем, сжатие только на уровне Journal-core. & Если для импликации требуется скрытая исключённая зависимость или предпосылки допустимости недостаточны, Source Theorem 3 должен быть сужен или отозван. \\
\occode{OBLIGATION::LEMMA_1} & Показать, что Source Theorem 3 закрывает Lemma 1 без расширения границы требования. & Поддержка исключительно через интерпретацию, поддержка исключительно через сжатие. & Если для Lemma 1 нужна опора, не прослеживаемая к Source Theorem 3, стержень разрывается и положительное ядро должно быть перестроено. \\
\occode{OBLIGATION::SOURCE_COROLLARY_3_2} & Показать, что Axiom 3.3 вместе с Definition 12.6 строго влечёт Source Corollary 3.2 и сохраняет необратимость. & Восстановление идентичности после законной смерти, метафорический язык "возрождения". & Если один и тот же континуум может вернуться численно идентичным после законной смерти без нарушения Axiom 3.3, следствие о необратимости должно быть пересмотрено или отозвано. \\
\occode{OBLIGATION::LEMMA_3} & Показать, что Source Corollary 3.2 закрывает Lemma 3 при объявленных границах остатка и возрождения. & Персистентность по аналогии, не объявленная эквивалентность возрождения. & Если остаток сохраняет полную живую идентичность, а не только структурный след, граница классификации должна быть переписана. \\
\occode{OBLIGATION::LEMMA_2} & Показать, что Lemma 1 вместе с Definition 12.5 влечёт Lemma 2 без скрытых расширений. & Ветви EXTENSION, пакеты BRIDGE_ONLY. & Если Lemma 2 требует исключённую ветвь или evidence BRIDGE_ONLY в качестве доказательной опоры, теоремный маршрут становится недействительным. \\
\occode{OBLIGATION::THEOREM_A} & Показать, что Lemma 2 достаточно для Theorem A и что извлечение journal-core не выходит за рамки master-теоремы. & Укрепление только на уровне Journal-core, исключённое семейство lift-heavy теорем. & Если ядро журнала требует опоры, отсутствующей на master-маршруте, статью нужно сузить или сначала починить master. \\
\bottomrule
\end{tabular}
\end{center}

\subsection{Backlog враждебного review}
\begin{center}
\begin{tabular}{@{}L{0.18\textwidth}L{0.25\textwidth}L{0.11\textwidth}L{0.15\textwidth}L{0.23\textwidth}@{}}
\toprule
ID review & Вызов & Статус & Состояние решения & Критерий выхода \\
\midrule
\occode{HOSTILE::LAWFUL_COLLAPSE} & Коллапс должен связываться с потерей допустимой реализации, а не с размытым метафорическим описанием. & \occode{PASS} & \occode{PASS_LOCKED} & Формальный досье связывает притязания на коллапс с потерей допустимого состояния, логикой порогов и источниковыми свидетельствами без утечек в прозу. \\
\occode{HOSTILE::IRREVERSIBILITY} & Необратимость по Axiom 3.3 должна запрещать восстановление численной идентичности после законной смерти. & \occode{PASS} & \occode{PASS_LOCKED} & Королларий о необратимости формализован как досье с явными условиями контрпримера и без лазейки идентичности. \\
\occode{HOSTILE::RESIDUE_REBIRTH_BOUNDARY} & Остаток и возрождение должны оставаться управляемыми классификациями, а не нарративными ярлыками. & \occode{PASS} & \occode{PASS_LOCKED} & Остаток, возрождение и персистентность разделены теоремой-граничными условиями с объявленными фальсификаторами. \\
\occode{HOSTILE::EXCLUDED_LIFT_FIREWALL} & Ни одно исключённое теоремное семейство с тяжёлым lift не должно быть тайно внесено в положительный ограниченный аргумент. & \occode{PASS} & \occode{PASS_LOCKED} & Реестр исключений показывает, что запрещённые семейства теорем и ветви не входят в маршрут ядра доказательства. \\
\occode{HOSTILE::JOURNAL_CORE_ASYMMETRY} & Journal-core должен оставаться законным извлечением из master и никогда не расширять требование. & \occode{PASS} & \occode{PASS_LOCKED} & Досье извлечения journal-core подтверждает, что каждый публичный теоремный тезис выводится трассируемо из маршрута master. \\
\occode{HOSTILE::K11_K12_IRREDUCIBILITY} & K11 и K12 должны быть доказаны необходимыми сверх K10 или явно демонтированы вместе с перестроением графа зависимостей. & \occode{PASS} & \occode{PROVED_NECESSARY_BEYOND_K10} & Либо dossier о неразложимости доказывает, что K11 и K12 несут невырезаемую научную работу, либо верхние уровни демпируются, и все зависимые маршруты перестраиваются. \\
\bottomrule
\end{tabular}
\end{center}

\subsection{Таблица приоритетов закрытия}
\begin{center}
\begin{tabular}{@{}L{0.16\textwidth}L{0.09\textwidth}L{0.09\textwidth}L{0.10\textwidth}L{0.18\textwidth}L{0.30\textwidth}@{}}
\toprule
Домен & Формальный порядок & Порядок затрат & Выигрыш на стоимость & Фаза & Примечание по параллельности \\
\midrule
Mathematics & 0 & 0 & \occode{116.666667} & \occode{PHASE_1_STABILIZE_CLOSED_CORE} & Работает непрерывно как формальный якорь для всех последующих доменов. \\
Physics & 1 & 1 & \occode{64.8} & \occode{PHASE_3_PHYSICS_CLOSURE} & Первичный канал первого закрытия; более поздний домен не должен опережать его. \\
Chemistry & 2 & 2 & \occode{32.0} & \occode{PHASE_4_CHEMISTRY_CLOSURE} & Следующий формальный канал после физики; часть подготовки данных может стартовать раньше, но продвижение ждёт фиксации физической оболочки. \\
Biology & 3 & 4 & \occode{17.818182} & \occode{PHASE_5_BIOLOGY_CLOSURE} & Формальный порядок предшествует системам, но высокая стоимость измерений делает этот канал медленнее. \\
Systems / Civilizational projection & 4 & 3 & \occode{16.615385} & \occode{PHASE_6_SYSTEMS_CLOSURE} & Более дешёвая эмпирическая подготовка допускается, но финальная унификация ждёт формального биологического маршрута. \\
\bottomrule
\end{tabular}
\end{center}

\subsection{Реестр абляции минимальности и неразложимости}
\begin{center}
\begin{tabular}{@{}L{0.20\textwidth}L{0.14\textwidth}L{0.23\textwidth}L{0.35\textwidth}@{}}
\toprule
Компонент & Тип & Статус абляции & Эффект удаления \\
\midrule
\occode{ROOT_THEOREM_STACK} & \occode{KERNEL_COMPONENT} & \occode{NECESSARY_UNDER_CURRENT_PROOF_STACK} & Обрушивает продвигаемую доктрину противоречия, подъёма lift или коллапса. \\
\occode{K_LEVEL_PARTITION} & \occode{K_LEVEL_BOUNDARY} & \occode{FRONTIER_PENDING_STRICT_IRREDUCIBILITY_PROOF} & Размоет явную доктрину уровней K и ослабит текущий реестр границ. \\
\occode{EMPIRICAL_ANCHORING_DISCIPLINE} & \occode{PROMOTION_GUARD} & \occode{NECESSARY_FOR_NON_RUBBERIZING_PROMOTION} & Позволит позитивным внешним притязаниям существовать без наблюдаемых, маршрутов данных и воспроизводимых фальсификаторов. \\
\occode{STRICT_UNIQUENESS_OR_MINIMALITY} & \occode{META_THEOREM} & \occode{EXPLICIT_FRONTIER_RESIDUE_LOCALIZED} & Самый сильный текущий результат остаётся границей невозможности, а не строгим доказательством уникальности. \\
\bottomrule
\end{tabular}
\end{center}

\subsection{Матрица конкуренции альтернативных моделей}
\begin{center}
\begin{tabular}{@{}L{0.22\textwidth}L{0.16\textwidth}L{0.20\textwidth}L{0.34\textwidth}@{}}
\toprule
ID конкуренции & Класс & Статус оценки & Объём сравнения \\
\midrule
\occode{SAME_CLAIM_CLASS::META_MODEL} & \occode{SAME_CLAIM_CLASS_METAMODEL} & \occode{PENDING_COMPLEXITY_PENALIZED_SAME_CLAIM_CLASS_COMPARISON} & Весь класс требований на объяснительно-прогнозную доменную прослойку ядра \\
\occode{DOMAIN_BASELINE::PHYSICS} & \occode{DOMAIN_BASELINE} & \occode{DOMAIN_BASELINE_COMPARISON_PASS_ELIGIBLE} & Ограниченный теоремно-родной физический пакет трассированно закрыт, replay-закрыт и компараторно-чист по закреплённому официальному held-out маршруту. \\
\occode{DOMAIN_BASELINE::CHEMISTRY} & \occode{DOMAIN_BASELINE} & \occode{DOMAIN_BASELINE_COMPARISON_PASS_ELIGIBLE} & Ограниченный теоремно-родной химический пакет трассированно закрыт, replay-закрыт и компараторно-чист по закреплённому термохимическому, спектральному и кинетическому held-out маршруту. \\
\occode{DOMAIN_BASELINE::BIOLOGY} & \occode{DOMAIN_BASELINE} & \occode{DOMAIN_BASELINE_COMPARISON_PASS_ELIGIBLE} & Ограниченный теоремно-родной биологический пакет трассированно закрыт, replay-закрыт и компараторно-чист через закреплённые мульти-наборы held-out биологии. \\
\occode{DOMAIN_BASELINE::SYSTEMS_CIVILIZATIONAL_PROJECTION} & \occode{DOMAIN_BASELINE} & \occode{DOMAIN_BASELINE_COMPARISON_PASS_ELIGIBLE} & Ограниченный теоремно-родной системный пакет трассированно закрыт, replay-закрыт и компараторно-чист на закреплённом held-out маршруте макротраекторий и смен режима. \\
\bottomrule
\end{tabular}
\end{center}

\subsection{Ограниченный бенчмарк сжатия}
\begin{center}
\begin{tabular}{@{}L{0.24\textwidth}L{0.12\textwidth}L{0.13\textwidth}L{0.41\textwidth}@{}}
\toprule
Ось & Единицы сложности & Значение эффективности & Интерпретация \\
\midrule
\occode{THEOREM_SUPPORT_DESCRIPTION_LENGTH} & 17 & 0.294118 & Сколько валидированного покрытия доменов сейчас приходится на одно корневое утверждение теоремной поддержки. \\
\occode{EMPIRICAL_PACKET_COMPLEXITY_VS_RESIDUAL_PERFORMANCE} & 5 & 1.0 & Сколько валидированного покрытия переносится на один активный числовой пакет при текущем replay-пороге. \\
\bottomrule
\end{tabular}
\end{center}
